This review makes a legislative case, for readers who shape medical and
artificial intelligence law without robotics or coding experience, that Physical
AI Trial Bill H. R. 9510 (v5.0) should be enacted into Federal law. It is built
on a single mechanism, verification before generation, in which a software agent
proposes a clinical action and a ten-gate verification, validation, and
uncertainty quantification examination either accepts it, escalates it to a
qualified human, or blocks it before it can reach a patient. The argument is
organized as eight emotional pillars that legislative-advocacy research finds most
persuasive: compassion, fear of preventable harm, moral outrage, hope,
responsibility, protection of vulnerable people, trust, and urgency. Each pillar
pairs a human appeal with a credible, cited fact, so the story is memorable and
the evidence is actionable. The review draws on a documented engineering lineage,
including a surgical-humanoid assurance run that passed 172 of 172 automated tests
across a ten-gate suite, and on the published advocacy literature describing how
testimony, coalition building, and policy entrepreneurship move a bill through
markup and reconciliation. Throughout, three quiet questions recur: how did we
ever ask humans alone to carry this task, why did we trust a clinical-trial system
with so many chances for compounding human error, and why would human peer review
remain the only check on increasingly complex and voluminous artificial
intelligence outputs. The conclusion is a single call to action with a specific
instrument attached.
